\subsubsection*{Ordinary split primes and the short band}
\label{subsec:ordinary-split-length}
\hypertarget{source-unramified-split-length}{}

After removing the type and ramified primes, the auxiliary field is unramified at the remaining primes in this range. At split primes the signed jumps lead to a short exceptional band; its placement must be compatible with the previously selected top positions.

\begin{established}[Unramified split range and signed jump window]\label{hyp:unramified-split-row}
Let \(\ell\) be outside the prime divisors of \(Q_S\), the primes
\(3,q,d\), and the two type primes. Consider first the range
\(\ell\le\ell_*\) together with the exceptional primes above
\(\ell_*\). At a prime split in the distinguished quadratic field, each column \(u\) is paired with its conjugate row, the column degrees sum to \(f\), and every completion is unramified. Put \(z=\log b\) and \(a=\lfloor b+S_c+\eta_c+t+2\rfloor\). For each column \(u\), both signed crystalline maps have jumps in \([-a,a]\), and their respective targets obey
\[
 \max\left\{\operatorname{rk}_O H^1_f(E_u,T_{2,u}),
          \operatorname{rk}_O H^1_f(E_{\bar u},T_{2,\bar u})\right\}
               \le [E_u:\mathbb Q_\ell]m .
\]
\end{established}
\begin{proof}
The excluded list contains every prime dividing the conductor of the host. Let \(U\) index the places on the selected side and \(D_u=[E_u:\mathbb Q_\ell]\). Conjugation gives the other side, so
\[
 2\sum_{u\in U}D_u=[E:\mathbb Q]=2f,\qquad |U|\le f .
\]

Lemmas~\ref{lem:marked-row-column-lattices} and~\ref{lem:marked-signed-jumps} identify the maps and the labelled jumps on this cut. At a good point use the direct normalized family on a column and
the Tate dual family on its conjugate row. Their jumps are in
\([-a,a]\) at every embedding by the jump boxes for the common
twist. Taking rational subquotients of finite direct sums preserves
crystallinity and this bound. For \(W\) on the column and
\(W^\vee(1)\) on the row, the free ranks of integral finite
cohomology are at most \(D_um\). Indeed the respective target
jumps avoid zero.

Among them \(k-1\) are positive on the column
and arise from negative jumps of \(W\) under \(j\mapsto-j-1\)
on the row. Write \(h_\tau^-\) for the number of negative target jumps at \(\tau\), including multiplicities. The split finite-rank formula in Section~\ref{subsec:split-local-length} gives, on either side,
\[
 \operatorname{rk}_O H^1_f(E_u,T_{2,u})=\sum_{\tau\mid u}h_\tau^-
                        \le D_um .
\] The column and row surjections in the congruence determinant
are on lattices in these subquotients. \label{cor:unramified-split-length-1}
\end{proof}

\begin{proposition}[Uniform finite lifting over the unramified completion]\label{prop:unramified-split-finite-lift}
For a crystalline lattice surjection in the signed window
\([-a,a]\) of Lemma~\ref{hyp:unramified-split-row}, the cokernel on integral finite cohomology
is killed by \(\ell^{e_A}\), where
\[
 e_A=\left\lceil23\bigl(2+2a/\ell+\log_\ell(2a+1)\bigr)\right\rceil .
\]
The allowance is independent of the degree of the coefficient
extension and the number of embeddings of the unramified field.
\end{proposition}

\begin{proof}
We show why the exponent on these surjections can be taken to be
\[
 e_A=\left\lceil23\bigl(2+2a/\ell+\log_\ell(2a+1)\bigr)\right\rceil .
\]
The estimate applies to a surjection of crystalline lattices with
the indicated window over any of these unramified fields. Use an
admitted coefficient extension containing the local embeddings and
the first cyclotomic level as in Section~\ref{subsec:cut-wach-models}.
The invariance of normalized length there returns the result on the
row or column surjection to its original coefficient ring. Use the
integral Wach--Herr comparison over the unramified field, with
Wach parameter \(\pi\) and height width \(h_1=2a\).
The Frobenius coefficient extraction of
Section~\ref{subsec:ramified-lift}, here with \(X=\pi\), puts
\(\pi^{h_1}N_2\) in the image of \(N_1\). On the full Wach
lattice the cyclotomic action is trivial modulo \(\pi\).
Choose a cyclotomic operator with character \(v\) generating
\(\mathbb Z_\ell^\times\), so its reduction has order
\(\ell-1\) and
\[
 v_\ell(v^{\ell-1}-1)=1.
\]

It acts on degree \(j\) by \(v^j\). Use the two interpolation
intervals \([0,h_1-1]\) and \([h_1,2h_1-1]\).
For a root \(v^i\) the only other roots giving nonunit
differences have indices congruent to \(i\) modulo \(\ell-1\).
On the two sides of \(i\) their numbers are
\[
 d_1=\lfloor i/(\ell-1)\rfloor,\qquad
 d_2=\lfloor(2h_1-1-i)/(\ell-1)\rfloor.
\]
The total valuation of the interpolation denominators satisfies
\[
 d_1+d_2+v_\ell(d_1!d_2!)\le \frac{2h_1\ell}{(\ell-1)^2}.
\]
The integral polynomial identity from the two intervals therefore puts \(\ell^{E_W}N_2\) in the image, for
\(E_W=\lfloor2h_1\ell/(\ell-1)^2\rfloor\).
For \(a=0\) the parameter-power inclusion already gives the lift.
Averaging over the subgroup \(\Delta\) of order \(\ell-1\)
preserves this bound on the invariant lattices. This elementary
interpolation bound suffices at the split caller.

Write \(N=N_0^\Delta\) for the invariant Wach lattice of the
kernel. The invariant coordinate \(P\) is a unit times
\(\pi^{\ell-1}\). Put \(Q'=P+\ell\), and let \(u\) be the
character of the principal cyclotomic generator \(\gamma\), with
\(v_\ell(u-1)=1\). The comparison gives the restricted coordinates
\((Q')^{-a}N\oplus PN\) for pairs and differential
\[
 d(y,z')=(\gamma-1)y-(\varphi-1)z'.
\]

It gives the following two facts on the kernel. First,
\(N\) is \(\ell\)-saturated in \((Q')^{-a}N\), and
\(\partial=(\gamma-1)/P\) preserves \(N\).
Second, if \(U\) is the set of source jumps, put
\[
 J=\big((U+\{1,\ldots,a\})\setminus\{0\}\big)
                \cup\{1,\ldots,a\},\qquad
 R_0(Z)=\prod_{j\in J}(Z-c_j),\quad c_j=(1-u^{-j})/\ell .
\]
Then \(\partial R_0(\partial)y\in N\) for \(y\in(Q')^{-a}N\).
The latter is the cyclotomic semisimplicity input on the
bounded-pole quotient. The comparison at the cyclotomic roots
reads the jumps shifted by the pole degrees \(1,\ldots,a\).
In the divided operator the eigenvalue for shift \(j\) is
\(c_j\), since the invariant coordinate at the root is
\(-\ell\). The extra factor \(Z\) accounts for shift zero.
Strictness puts the kernel jumps among the source jumps.
\label{cor:unramified-split-length-2}

Here is the correction of a defect in lifting a pair. Put
\(d_0=\lfloor a/(\ell-1)\rfloor\) and
\(G=\prod_{i=1}^{d_0}(1-u^{(\ell-1)i})\). On
\(P^iN/P^{i+1}N\), \(\gamma\) acts by \(u^{(\ell-1)i}\).
Successive subtraction through the graded factors gives
\[
 GPN\subseteq(\gamma-1)PN+P^{d_0+1}N.
\]

The height inclusion and the factor \((Q')^{\ell-1}\) in
\(\varphi(P)/P\) make Frobenius contracting on
\(P^{d_0+1}N\), since \((\ell-1)(d_0+1)>a\).
In this complete submodule the geometric series solves the inhomogeneous equation:
\[
 (1-\varphi)\sum_{n\ge0}\varphi^n x=x\qquad(x\in P^{d_0+1}N).
\] The integral restricted Herr
correction applied with these data kills every rationally
trivial defect by \(G R_0(0)\). To see the role of division
by \(P\), write the defect as \(P\beta\).

In the equation
for a multiple which is a boundary, use the graded correction
on the second coordinate. Dividing by \(P\) leaves
\(\partial\) on the first coordinate and an integral term
on the other side. Applying \(R_0(\partial)\) makes this
first-coordinate difference integral, by the bounded-pole
identity. More explicitly this gives \(\ell^j G R_0(\partial)\beta\in N\)
when the original boundary equation holds after multiplication by \(\ell^j\). Saturation in the
bounded-pole lattice and \(G\mid\ell^{v_\ell(G)}\) up to a unit
give
\[
 R_0(\partial)\beta\in N.
\]

Using \(R_0(Z)=R_0(0)+Z S_0(Z)\) and
\(P\partial=\gamma-1\), the corrected defect satisfies
\[
  R_0(0)P\beta
       =P R_0(\partial)\beta-(\gamma-1)S_0(\partial)\beta .
\]
Here \(P R_0(\partial)\beta\in PN\). The second summand is
the indicated cyclotomic difference. The positive-degree
correction by \(G\) and the deep correction by \(1-\varphi\)
apply to the first summand in \(PN\).

\label{cor:unramified-split-length-3}

An integral finite class is represented by a crystalline extension
of the trivial lattice. Lifting its quotient generator on the
invariant Wach module costs \(E_W\) and gives a restricted pair.
Lifting the two coordinates through the given surjection costs
a second \(E_W\). The defect lies in the kernel coordinates
by exactness. It is rationally trivial by surjectivity on rational
finite classes for crystalline quotients and the rational Herr
comparison. The integral correction gives a finite source class.
The resulting exponent is at most
\[
  2E_W+\sum_{j\in J}v_\ell(j)
       +(d_0+1)+v_\ell((d_0+1)!).
\]
Indeed \(v_\ell(c_j)=v_\ell(j)\) and
\(v_\ell(G)=d_0+v_\ell(d_0!)\).
The distinct indices lie in \([-a,2a]\setminus\{0\}\), so their
sum of valuations is at most \(3a/(\ell-1)\).

For \(\ell\ge3\), the three displayed contributions are at most
\(9h_1/\ell\), \(9h_1/(4\ell)\), and
\(3(a/(\ell-1)+1)/2\), respectively. Their sum is at most
\(13h_1/\ell+2\le e_A\). In particular the exponent has no
factor depending on the number of embeddings or the coefficient
ramification. \label{cor:unramified-split-length-4}
\end{proof}

\paragraph{The determinant on the split family.}
\begin{lemma}[Split determinant bound at the same lift]
\label{lem:unramified-split-determinant-total}
Use the unramified split range and the maps of
Lemma~\ref{hyp:unramified-split-row}, with \(z\ge1\). If \(L\) denotes the total lifting term and \(D\) the sum of the remaining local determinant terms, then
\[
 L\le (2fm+f)e_A,\qquad D\le7200fnt^2Rz^2.
\]
The separate global and signature allowances are unchanged.
For the inert declaration only \(B_0\) from this calculation
is used.
\end{lemma}

\begin{proof}
Use Proposition~\ref{prop:unramified-split-finite-lift} on the
row and column surjections specified by the same lift.
Normalize every length by \(\operatorname{length}_O(O/\ell O)\).
Write \(R=r+t\), \(h=\lfloor2b\rfloor\),
\(N_0=4nR\), \(K_0'=\lfloor\log_\ell(2h)\rfloor\) and
\(K_1'=\lfloor\log_\ell(2(h+1))\rfloor\).
There are at most \(D_wR\le N_0\) distinct jumps in the
character calculation. Passing to \(E_w(\mu_\ell)\) has
prime-to-\(\ell\) degree. Thus the character calculation giving
\eqref{eq:tag-r2} applies, with depth at most \(4+K_0'\), and gives exponent
\[
 c_B=6+K_0'+7t+tN_0(4+2K_0').
\]
For the column and dual row torsion comparison the window lies in
\([-h-1,h+1]\) with zero absent.

In the invariant formula in
Section~\ref{subsec:split-local-length}, take one truncation
degree and full cyclotomic depth one. The depth term is at most
\(3+K_1'\), the sum over jumps is at most \(N_0(4+2K_1')\),
and the two integral comparison losses in this degree add two.
The bound for the invariant-killing exponent is therefore
\[
 c_0=5+K_1'+N_0(4+2K_1').
\]
If the jump sum is empty, the extra unit allowance in the
maximum is also covered by this bound since \(N_0\ge1\).
Consequently \(h_W(u)\le rc_0\).
The free and torsion parts of the dual row give colength at
most \(D_um e_A+rc_0\). The torsion bound follows from the
invariant subgroup of the divisible target modulo its lattice.
These are estimates on the given lattices. \label{cor:unramified-split-length-5}

Write \(k_c\) for the number of columns. At depth one the
marked scalar satisfies \(v_\Delta\le 2t\) by the cyclotomic
difference valuations. On the nontrivial scalar twists
\((K_\ell/O)(-j)\) with \(1\le j<t\), the unramified base has
an inertia element of character \(1+\ell\), so its invariants are killed by \((1+\ell)^{-j}-1\). Their
normalized length is bounded by
\[
 v_\ell((1+\ell)^{-j}-1)=v_\ell((1+\ell)^j-1)
                         =1+v_\ell(j)\le t .
\] Use \(h_0\le t\) for
their maximum, or zero for an empty set.

Summing the three
blocks of the determinant gives a common local term
\[
 B_0=t\sum_{w\ {\rm row}}(D_w+2)c_B
       +4R(k_c+1)v_\Delta+tk_ch_0 .
\]
The column and conjugate row at \(u\) add
\((2D_um+1)e_A\) and the error
\(D_umv_\Delta+2v_\Delta+h_W(u)+rc_0\).
In fact the column contributes \(j_u(e_A+v_\Delta)+e_A\)
and the row contributes \(D_um e_A\) before adding their
torsion costs. The other column and coboundary denominators
are those of Section~\ref{subsec:split-local-length}.

Using \(z\ge1\), \(K_0'\le2+2z\) and \(K_1'\le3+2z\), the two scalar losses are
\[
 c_B\le120t n\,R z,\qquad c_0\le120 n\,R z.
\]
Put
\(U_0=120fnt^2Rz\). The sum of row degrees plus twice their
number is at most \(6f\), and \(k_c\le2f\). Thus the three
terms of \(B_0\) are at most \(6U_0,24U_0,2U_0\).
The additional split errors are at most \(4U_0,8U_0,4U_0\)
for the degree term, the scalar term, and the two invariant
terms, using \(m\le r\le t\).

Adding the six errors bounds their sum by
\[
        (6+24+2+4+8+4)U_0
           \le 7200 fnt^2Rz^2 ,
\]
since \(z\ge1\). The sum of the lifting terms is at most
\((2fm+f)e_A\). Together these give the split and boundary
entries in Section~\ref{subsec:prime-cost-sums}, since
\(e_A\ll 1+H/\ell+\log_\ell H\). The global kernels, the
signature bound and the tame character costs retain their
separate allowances. In the inert case \(k_c=0\) and the
same calculation bounds \(B_0\) alone. The rank-\(r\) row
colengths there depend on the inert jump word.
\label{cor:unramified-split-length-6}
\end{proof}

\paragraph{Transport of the band at an inert place.}
\begin{lemma}[Passage from the row band to the place bound]
\label{lem:inert-place-band}
In the last two inert ranges let \(b_{\rm row}\) be the marked row band indicator and \(b_\ell\) the place indicator used in the weighted sum at its rational prime. Then
\[
 b_{\rm row}\le b_\ell\qquad(b_{\rm row},b_\ell\in\{0,1\}).
\]
\end{lemma}

\begin{proof}
On a row in either of the last two inert ranges, suppose the
local band indicator is one. In the canonical data this gives
a good point, a row word at a place over \(\ell\ge11\), and
a position with smaller jumps of local weight \(a_0\) with
\[
 |\ell a_0-b_0|\le K\{\ell\eta+\ell S+b_0/\ell+a_0+
                                          \ell\log_\ell(2h_b)\}.
\]
Here \(h_b\) is the height of the word, and \(b_0\) its large
weight. The weight box on the same word supplies a place weight
\(s\) between \(\omega\) and the upper place-weight size with
\[
 |a_0-s|\le T,\qquad b-T\le b_0\le b-T+1,\qquad T=(t+1)/2 .
\]
The prime beneath this row equals the table prime, as the
contraction of its maximal ideal contains a unique rational
prime.

We have \(\eta\le\eta_3=\eta_c+t+2\),
\(S\le S_c\), \(h_b\le b_0+\eta+S\le 2b\) on the
selected range, and \(\log_\ell(4b)\le z\) since \(z\ge1\)
and \(4e<11\). Put \(A_s=\ell(\eta_3+S_c+z)+b/\ell+s\).

The braces in the local inequality are at most
\(A_s+\ell\eta_3\), since \(a_0\le s+T\) and
\(b_0/\ell\le b/\ell+1/\ell\).
The two changes on the left cost at most
\(\ell T+T\le\ell\eta_3\). Hence
\[
                 |\ell s-b|\le(2K+1)A_s .
\]
The table constant dominates \(6K+1\), and therefore the
indicator in the weighted sum dominates every local band
indicator. Hence the global harmonic sum bounds the frequency of the band on each of these rows. Its contribution per occurrence still requires an inert-row colength estimate.
\label{cor:unramified-split-length-7}
\end{proof}

